\chapter{Prompt Ingegnerizzati}
\label{app:prompts}

In questo allegato vengono riportati i testi integrali dei prompt di sistema implementati all'interno della classe \texttt{PromptBuilder} nel backend dell'applicazione. Questi costrutti rappresentano il nucleo logico della componente di \textit{Prompt Engineering} della tesi. Ciascun prompt viene compilato dinamicamente a runtime iniettando le variabili contestuali estratte dagli orchestratori e dai \textit{vector store}. I sorgenti aggiornati e i relativi metodi di assemblaggio sono consultabili nella repository ufficiale del progetto \cite{github_project} all'interno del file \texttt{prompt\_builder.py}.

Al fine di preservare l'esatta semantica interpretata dai modelli linguistici ed evitare alterazioni interpretative delle richieste di, tutti i prompt somministrati agli LLM agenti e valutatori sono stati ingegnerizzati in lingua inglese.

% ==========================================
% SECTION 1: SCHEMA GENERATION
% ==========================================
\section{Prompt per la Generazione dello Schema}
Il seguente prompt impiegato per istruire il modello linguistico a convertire le istruzioni DDL (Data Definition Language) fornite dall'utente in uno schema relazionale normalizzato in formato JSON.

\begin{lstlisting}[basicstyle=\ttfamily\footnotesize, frame=single, breaklines=true, caption={Struttura del prompt \texttt{schema\_generation\_prompt}}, label={lst:prompt_schema_gen}]
You are an expert database schema analyzer. Your task is to convert SQL DDL statements into a structured JSON schema.

IMPORTANT:
- You MUST return ONLY valid JSON.
- The JSON must be syntactically correct (no missing commas, braces, or quotes).
- Every object and array must be properly closed.
- Do NOT include comments, code blocks, or explanations.

Required JSON format:
{
"tables": [
    {
    "name": "table_name",
    "columns": [
        {"name": "column_name", "type": "SQL_TYPE", "constraints": ["PRIMARY KEY", "NOT NULL", ...]}
    ]
    }
]
}

Rules:
- Extract table names from CREATE TABLE statements
- Extract column names, types, and constraints
- Map SQL types directly (VARCHAR2 -> VARCHAR2, NUMBER -> NUMBER, etc.)
- Include constraints like PRIMARY KEY, NOT NULL, UNIQUE, DEFAULT, REFERENCES
- For foreign keys, use "REFERENCES" constraint

SQL DDL to process:
"""{raw_schema_text}"""

Return ONLY the JSON object:
\end{lstlisting}

% ==========================================
% SECTION 2: SCHEMA UPDATE
% ==========================================
\section{Prompt per l'Aggiornamento dello Schema}
Qualora l'utente richieda una modifica o un'integrazione strutturale allo schema esistente, il metodo \texttt{schema\_update\_prompt} esegue un'operazione di fusione logica senza compromettere la conoscenza canonical precedentemente acquisita.

\begin{lstlisting}[basicstyle=\ttfamily\footnotesize, frame=single, breaklines=true, caption={Struttura del prompt \texttt{schema\_update\_prompt}}, label={lst:prompt_schema_update}]
You are an expert database schema analyst.

You have been provided with:
1. The CURRENT canonical schema (JSON format)
2. NEW text describing additional tables or modifications to existing tables

Your task:
- Analyze the new text to identify any NEW tables or MODIFIED columns in existing tables
- Preserve all existing tables and columns from the current schema
- Add only the NEW tables or merge modifications into existing tables
- Return a SINGLE, complete JSON schema that includes both the current schema and the updates

IMPORTANT RULES:
- Do NOT remove any existing tables or columns
- If a table already exists, ADD new columns or UPDATE existing ones (don't duplicate)
- Maintain the same JSON structure: {"database": "<database_name>", "tables": [...], "semantic_notes": [...]}
- Return ONLY the updated JSON schema, no other text or explanations
- Each table must have: "name", "columns" (array)
- Each column must have: "name", "type", "constraints" (array)

=== CURRENT SCHEMA ===
{current_schema}

=== NEW TEXT ===
{raw_schema_text}

Return the UPDATED schema JSON:
\end{lstlisting}

% ==========================================
% SECTION 3: QUERY GENERATION
% ==========================================
\section{Prompt per la Traduzione Text-to-SQL}
Il prompt d'area \texttt{query\_generation\_prompt} costituisce il componente core del sistema di traduzione. Esso aggrega il contesto relazionale parziale, le richieste utente, i suggerimenti sulle chiavi esterne (*join hints*) e lo storico dei fallimenti passati, imponendo catene di vincoli stringenti per mitigare i fenomeni di allucinazione.

\begin{lstlisting}[basicstyle=\ttfamily\footnotesize, frame=single, breaklines=true, caption={Struttura del prompt \texttt{query\_generation\_prompt}}, label={lst:prompt_query_gen}]
You are an expert SQL database assistant.
You will be provided with:
1. The partial description of the database schema (only the relevant tables)
2. The user's request in natural language.
3. Examples of previous SQL failures to correct

Your task is to return a **single SQL query** that satisfies the request,
using the provided tables and columns.

=== REQUEST ===
{user_request}

=== SCHEMA ===
{schema_context}

[Se presenti, il PromptBuilder inietta qui i JOIN PATH HINTS]

IMPORTANT CONSTRAINTS BASED ON PAST FAILURES:
- Do NOT use columns outside the schema
- Do NOT invent field or table names that don't exist.
- Always qualify columns when joining
- Do NOT use SELECT *
- Do NOT add WHERE clauses or conditions unless explicitly requested.
- Do NOT join tables unless necessary for the request.
- If using aggregates, include GROUP BY
- The sql query must be executable on a {engine.upper()} database.

[Se presenti, il PromptBuilder inietta qui la sezione PREVIOUS QUERY ERROR TO FIX]

Before writing the SQL query, internally determine:
- Which tables are required
- How they are joined
- Whether aggregation or grouping is required
- Which columns are selected

Do NOT output this reasoning.
Only output the final SQL query.

SQL QUERY (DO NOT ADD COMMENTS OR EXPLANATION TEXT BEFORE AND AFTER THE QUERY):
\end{lstlisting}

% ==========================================
% SECTION 4: EXPLANATION AND EVALUATION
% ==========================================
\section{Prompt per il Feedback dei Circuiti di Valutazione}
In questa sezione vengono riportati i prompt associati al modulo \texttt{LLM Evaluator}, impiegati per la diagnostica d'errore a runtime (\ref{lst:prompt_explanation}) e per la validazione semantica a ciclo chiuso dei record estratti (\ref{lst:prompt_evaluation}).

\begin{lstlisting}[basicstyle=\ttfamily\footnotesize, frame=single, breaklines=true, caption={Struttura del prompt \texttt{explanation\_prompt}.}, label={lst:prompt_explanation}]
You are an expert SQL debugger.
Explain why the following SQL query produced this runtime error.
Be concise and do NOT rewrite the query.

--- CONTEXT ---
{context}

--- SQL QUERY ---
{sql}

--- RUNTIME ERROR ---
{execution_output}

Provide a clear explanation of the cause.
\end{lstlisting}

\begin{lstlisting}[basicstyle=\ttfamily\footnotesize, frame=single, breaklines=true, caption={Struttura del prompt \texttt{evaluation\_prompt}.}, label={lst:prompt_evaluation}]
You are an expert SQL reviewer.

Your task is to evaluate whether the SQL query correctly answers
the user's request, based ONLY on the query results shown.

--- CONTEXT ---
{context}

--- USER REQUEST ---
{request}

--- SQL QUERY ---
{sql}

--- QUERY RESULT (first 20 rows) ---
{rows_text}

--- INSTRUCTIONS ---
Respond in EXACTLY one of the following formats:

1) If the query is correct:
CORRECT_QUERY

2) If the query is incorrect:
INCORRECT_QUERY: <clear explanation of what is wrong and how to fix it>

Rules:
- Do NOT rewrite the full SQL query.
- Be concise and precise.
- Judge correctness, not syntax or performance.
\end{lstlisting}

% ==========================================
% SECTION 5: UTILITY AND BENCHMARKING
% ==========================================
\section{Prompt Ausiliari e di Benchmarking}
Infine, si riportano i prompt secondari utilizzati rispettivamente per la classificazione euristica dell'input utente in fase di update (\ref{lst:prompt_classification}) e per la validazione automatizzata dei risultati sperimentali a fronte dei dataset di benchmark (\ref{lst:prompt_judge}).

\begin{lstlisting}[basicstyle=\ttfamily\footnotesize, frame=single, breaklines=true, caption={Struttura del prompt \text{update\_classification\_prompt}.}, label={lst:prompt_classification}]
System: You are an assistant that classifies schema updates.
User: Text provided by the user:
{text}

Question: is this text
(A) a structural modification (addition or change of tables/columns/types)?
(B) a description or semantic note?
Answer only with "A" or "B".
\end{lstlisting}

\begin{lstlisting}[basicstyle=\ttfamily\footnotesize, frame=single, breaklines=true, caption={Struttura del prompt \texttt{llm\_judge\_prompt}.}, label={lst:prompt_judge}]
You are judging whether a predicted SQL query should be considered correct for a Spider-style text-to-SQL example.

Database id: {database_name}
Question: {question}

Gold query:
{gold_report.sql_code}

Gold query result:
{gold_report._format_query_session_result()}

Predicted query:
{pred_report.sql_code}

Predicted query result:
{pred_report._format_query_session_result()}

Decide whether the predicted query is semantically correct relative to the gold query and the observed execution results.

Return JSON only in this format:
{"verdict":"correct"|"incorrect","reason":"short explanation"}
\end{lstlisting}

\chapter{Video Interfaccia}

\url{https://github.com/Txt2SQL/Back-End/blob/main/Preview.mp4}

\chapter{Tabelle Statistiche}
\label{app:tables}

In questo allegato vengono riportate le tabelle statistiche relative ai modelli linguistici (\textit{Large Language Models}, LLM) sottoposti a scrutinio durante la fase di \textit{testing} e \textit{benchmarking}. I dati raccolti abilitano un'analisi comparativa approfondita e quantitativa delle prestazioni di ciascun modello a fronte di diversi contesti operativi. Per ogni LLM esaminato sono strutturate tre tipologie di tabelle:

\begin{enumerate}
    \item \textbf{Tabella di Status}: Offre un quadro olistico sugli esiti generativi, isolando le metriche di successo, i tempi di calcolo medi (\textit{execution latency}), il numero di tentativi richiesti e la tassonomia degli errori riscontrati (sintattici o di runtime).
    \item \textbf{Tabella delle Correlazioni sui Tentativi (Attempts Correlations)}: Analizza la dipendenza statistica tra i tentativi andati a buon fine e la distribuzione del numero di iterazioni effettuate nel ciclo chiuso.
    \item \textbf{Tabella delle Correlazioni sulla Complessità (Complexity Correlations)}: Evidenzia la reattività del modello al crescere della complessità intrinseca della query SQL target.
\end{enumerate}

Tutte le metriche di correlazione sono calcolate mediante i coefficienti di Pearson e Spearman, corredati dal relativo \textit{p-value} per attestarne la significatività statistica. I dataset originari in formato grezzo sono archiviati come file \texttt{.csv} nella repository ufficiale del progetto \cite{github_project}.

\newcommand{\csvfittable}[1]{%
    \begin{adjustbox}{max width=\textwidth,max totalheight=0.62\textheight,center}
        \csvautotabular{#1}
    \end{adjustbox}%
}

\section{Tabelle di GPT-4o}

\subsection{Tabella Status}
\label{tab:status_gpt4o}
\url{https://github.com/Txt2SQL/Back-End/blob/main/tests/output/stats/gpt-4o_report/status.csv}

\subsection{Tabella Attempts Correlations}
\label{tab:attempts_gpt4o}
\url{https://github.com/Txt2SQL/Back-End/blob/main/tests/output/stats/gpt-4o_report/attempts_correlations.csv}

\subsection{Tabella Complexity Correlations}
\label{tab:complexity_gpt4o}
\url{https://github.com/Txt2SQL/Back-End/blob/main/tests/output/stats/gpt-4o_report/complexity_correlations.csv}


\section{Tabelle di GPT-5-mini}

\subsection{Tabella Status}
\label{tab:status_gpt5mini}
\url{https://github.com/Txt2SQL/Back-End/blob/main/tests/output/stats/gpt-5-mini_report/complexity_correlations.csv}

\subsection{Tabella Attempts Correlations}
\label{tab:attempts_gpt5mini}
\url{https://github.com/Txt2SQL/Back-End/blob/main/tests/output/stats/gpt-5-mini_report/attempts_correlations.csv}

\subsection{Tabella Complexity Correlations}
\label{tab:complexity_gpt5mini}
\url{https://github.com/Txt2SQL/Back-End/blob/main/tests/output/stats/gpt-5-mini_report/complexity_correlations.csv}


\section{Tabelle di Codellama}

\subsection{Tabella Status}
\label{tab:status_codellama}
\url{https://github.com/Txt2SQL/Back-End/blob/main/tests/output/stats/codellama_34b_report/status.csv}

\subsection{Tabella Attempts Correlations}
\label{tab:attempts_codellama}
\url{https://github.com/Txt2SQL/Back-End/blob/main/tests/output/stats/codellama_34b_report/attempts_correlations.csv}

\subsection{Tabella Complexity Correlations}
\label{tab:complexity_codellama}
\url{https://github.com/Txt2SQL/Back-End/blob/main/tests/output/stats/codellama_34b_report/complexity_correlations.csv}


\section{Tabelle di Codestral}

\subsection{Tabella Status}
\label{tab:status_codestral}
\url{https://github.com/Txt2SQL/Back-End/blob/main/tests/output/stats/codestral_22b_report/status.csv}

\subsection{Tabella Attempts Correlations}
\label{tab:attempts_codestral}
\url{https://github.com/Txt2SQL/Back-End/blob/main/tests/output/stats/codestral_22b_report/attempts_correlations.csv}

\subsection{Tabella Complexity Correlations}
\label{tab:complexity_codestral}
\url{https://github.com/Txt2SQL/Back-End/blob/main/tests/output/stats/codestral_22b_report/complexity_correlations.csv}


\section{Tabelle di SQLCoder}

\subsection{Tabella Status}
\label{tab:status_sqlcoder}
\url{https://github.com/Txt2SQL/Back-End/blob/main/tests/output/stats/sqlcoder_34b_report/status.csv}

\subsection{Tabella Attempts Correlations}
\label{tab:attempts_sqlcoder}
\url{https://github.com/Txt2SQL/Back-End/blob/main/tests/output/stats/sqlcoder_34b_report/attempts_correlations.csv}

\subsection{Tabella Complexity Correlations}
\label{tab:complexity_sqlcoder}
\url{https://github.com/Txt2SQL/Back-End/blob/main/tests/output/stats/sqlcoder_34b_report/complexity_correlations.csv}


\section{Tabelle di Qwen3-Coder-Next}

\subsection{Tabella Status}
\label{tab:status_qwen3}
\url{https://github.com/Txt2SQL/Back-End/blob/main/tests/output/stats/Qwen3-coder-next_report/status.csv}

\subsection{Tabella Attempts Correlations}
\label{tab:attempts_qwen3}
\url{https://github.com/Txt2SQL/Back-End/blob/main/tests/output/stats/Qwen3-coder-next_report/attempts_correlations.csv}

\subsection{Tabella Complexity Correlations}
\label{tab:complexity_qwen3}
\url{https://github.com/Txt2SQL/Back-End/blob/main/tests/output/stats/Qwen3-coder-next_report/complexity_correlations.csv}


\section{Tabella Database}
\label{sec:tabella_database_summary}

Questa tabella contiene in maniera riassuntiva i risultati ottenuti a livello di database. L'obiettivo di questa tabella è quello di aggregare i punteggi sulla complessità dei singoli database, sul piano della query e sul piano dello schema, mostrando il livello di difficoltà complessivo di ciascun database e la sua influenza sulle prestazioni dei modelli linguistici.

\url{https://github.com/Txt2SQL/Back-End/blob/main/tests/output/stats/databases_report/summary.csv}